\chapter{Background and Related Work}
\label{chap:background}

This chapter reviews the concepts needed for the later experiments: spiking-neuron models, information-theoretic and dynamical notation, reservoir computing, graph-based EEG analysis, affective task labels, and interpretability. The review also states where cited theory does not transfer directly to the finite software implementation or the SHAPE data.

%% ===================================================================
%% SECTION 1: NEURAL COMPUTATION PARADIGMS
%% ===================================================================
\section{The Evolution of Neural Computation Paradigms}
\label{sec:neural_paradigms}

\subsection{From Binary Logic to Rate-Based Abstraction}
McCulloch and Pitts formalized a neuron as a binary threshold unit \cite{mcculloch_logical_1943}, and Rosenblatt's perceptron learned a linear decision rule from data \cite{rosenblatt_perceptron_1958}. A single perceptron is limited to linearly separable problems. Later networks used continuous activations and multiple layers to represent broader classes of functions.

Continuous-activation networks are commonly trained by backpropagation and can approximate broad function classes under stated architectural assumptions. Their activations are sometimes interpreted as rate-like abstractions, although artificial units are not literal biological firing-rate models. Rapid visual processing has motivated neural codes that use spike timing as well as rate \cite{simon_j_thorpe_biological_1989}; it does not establish that rate information is absent or that one coding account explains all cortical computation.

\subsection{Spiking Neural Networks}
Spiking Neural Networks represent state with neuron variables and communicate through discrete events. Their implementations may be event driven or advanced on a simulation clock, as in this dissertation. Under a particular formal comparison, Maass showed that spiking networks can realize some functions with fewer units than specified sigmoidal networks \cite{maass_networks_1997}. That result is not a general guarantee of accuracy, biological fidelity, or energy efficiency for the reservoir used here.

%% ===================================================================
%% SECTION 2: SPIKING NEURON MODELS
%% ===================================================================
\section{The Biophysics and Mathematics of Spiking Neuron Models}
\label{sec:neuron_models}
The neuron model determines the state update and spike rule. Its selection trades biological detail against computational cost and parameter identifiability.

\subsection{The Leaky Integrate-and-Fire (LIF) Model}
A canonical and widely used model is the Leaky Integrate-and-Fire (LIF) neuron \cite{gerstner_spiking_2002}. It abstracts the neuron's membrane as a simple RC circuit. The subthreshold dynamics of the membrane potential $v(t)$ are described by the first-order linear ordinary differential equation:
\begin{equation}
    \tau_m \frac{dv(t)}{dt} = -(v(t) - v_{\text{rest}}) + R_m I(t)
    \label{eq:lif_continuous}
\end{equation}
where $\tau_m = R_m C_m$ is the membrane time constant, $v_{\text{rest}}$ is the resting potential, $R_m$ is the membrane resistance, and $I(t)$ is the total input current. When $v(t)$ reaches a threshold $v_{\text{th}}$, two events occur: a spike is emitted to downstream neurons, and the potential is reset to a value $v_{\text{reset}}$ for a brief refractory period during which it cannot spike. The LIF model's primary limitation is its inability to reproduce the rich diversity of firing patterns (e.g., bursting, adaptation) seen in biological neurons.

\subsection{Beyond the LIF: The Izhikevich and AdEx Models}
To capture this richer dynamical repertoire, more complex models have been proposed. The Izhikevich model, for instance, uses a two-dimensional system of ODEs to couple the membrane potential $v$ with a slower membrane recovery variable $u$:
\begin{align}
    \frac{dv}{dt} &= 0.04v^2 + 5v + 140 - u + I \\
    \frac{du}{dt} &= a(bv - u)
\end{align}
with a reset condition: if $v \ge 30$\,mV, then $v \leftarrow c,\; u \leftarrow u+d$. Different parameter settings reproduce a range of canonical firing patterns \cite{izhikevich_simple_2003}. The Adaptive Exponential Integrate-and-Fire model is another two-dimensional alternative that represents adaptation and nonlinear spike initiation.

\subsection{Justification for the LIF Model in this Dissertation}
\label{sec:lif_justification}
The LIF model was selected because its state and reset are simple to inspect and inexpensive to simulate at the reservoir sizes tested here. This choice does not show that richer neuron dynamics are unimportant; comparing neuron models is outside the empirical scope of the dissertation.

The subthreshold response of a single LIF unit is exponential. Lindeberg derives time-causal scale-space kernels under a particular set of axioms \cite{lindeberg_2016}, which provides mathematical context for exponential temporal smoothing. The recurrent, thresholded, and reset reservoir is not itself that linear scale-space system.

\subsection{Discrete-Time Implementation and the \texorpdfstring{$\beta$--$\tau_m$}{beta--tau-m} Relationship}
\label{sec:discrete_lif}
The experiments use a custom NumPy implementation of a discrete-time LIF-like reservoir; they do not call \texttt{snnTorch}'s \texttt{RLeaky} cell. Before thresholding, the membrane update for neuron $i$ is
\begin{equation}
    M_i[t] = (1 - \beta)\, M_i[t-1]\,(1 - S_i[t-1]) + I_{\text{tot},i}[t]
    \label{eq:lif_discrete}
\end{equation}
where $\beta\in(0,1)$ is the leak fraction, $S_i[t-1]\in\{0,1\}$ is the previous spike indicator, and $I_{\text{tot},i}[t]$ is the current input plus recurrent drive. The factor $(1-S_i[t-1])$ clears the retained membrane after a prior spike. The implementation then sets $S_i[t]=\mathbf{1}[M_i[t]\geq M_{\text{th}}]$, subtracts $S_i[t]M_{\text{th}}$, and clips the membrane below at zero. This combination is the exact update implemented in the archived scripts; it is not identical to the default reset semantics of every LIF library.

For a forward-Euler approximation to Equation~\eqref{eq:lif_continuous}, the retained fraction is $1-\Delta t/\tau_m$. Matching it to $1-\beta$ gives
\begin{equation}
    \beta = \frac{\Delta t}{\tau_m}, \qquad \tau_m = \frac{\Delta t}{\beta}.
    \label{eq:beta_tau}
\end{equation}
Equation~\eqref{eq:beta_tau} is exact only within that Euler discretization. Matching the exact exponential solution instead gives $1-\beta=e^{-\Delta t/\tau_m}$, or $\tau_m=-\Delta t/\ln(1-\beta)$. With $\beta=0.05$, these mappings give $20$ and approximately $19.5$ simulation steps, respectively. Any conversion to milliseconds must also use the actual post-decimation step size; the parameter is not itself a measured biological membrane constant.


%% ===================================================================
%% SECTION 3: MATHEMATICAL PRELIMINARIES
%% ===================================================================
\section{Mathematical Preliminaries}
\label{sec:math_prelim}

This section introduces the notation and theoretical results used in the later chapters. Each result is stated only to the extent needed to interpret the implemented analyses.

\subsection{Information-Theoretic Foundations}
\label{sec:info_theory}

Several constructs from information theory are used in this dissertation, particularly in the feature-space analysis of Chapter~4 and the dynamical characterization framework of Chapter~6.

\paragraph{Mutual Information.}
For two random variables $X$ and $Y$ with joint distribution $p(x,y)$ and marginals $p(x)$, $p(y)$, the mutual information is defined as
\begin{equation}
    I(X; Y) = \sum_{x,y} p(x,y) \log_2 \frac{p(x,y)}{p(x)\,p(y)}\,
    \label{eq:mutual_info}
\end{equation}
with the continuous analogue replacing sums by integrals \cite{cover_elements_2006}. Mutual information quantifies the reduction in uncertainty about $Y$ obtained by observing $X$ (and vice versa). It is non-negative, symmetric, and zero if and only if $X$ and $Y$ are statistically independent.

\paragraph{The Data Processing Inequality.}
If $X \to Y \to Z$ forms a Markov chain, meaning $Z$ is conditionally independent of $X$ given $Y$, then
\begin{equation}
    I(X; Z) \leq I(X; Y).
    \label{eq:dpi}
\end{equation}
This inequality is fundamental to the interpretation of reservoir-derived features: any processing of the EEG signal through the reservoir and subsequent feature extraction can only preserve or lose information about the latent state that generated the signal; it cannot create new information \cite{cover_elements_2006}. When Chapter~4 reports that the Fisher Discriminant Ratio increases after the reservoir transformation, the correct interpretation is that the reservoir has improved the \emph{representation geometry} for the chosen decoder, not that it has increased the mutual information between the signal and the class label. Chapter~6 invokes the DPI explicitly as the foundational constraint bounding what reservoir-derived dynamical metrics can claim about the input process.

\paragraph{Fisher Discriminant Ratio.}
For a two-class problem with class-conditional mean vectors $\boldsymbol{\mu}_0$, $\boldsymbol{\mu}_1$ and within-class scatter matrix $\mathbf{S}_W = \sum_{i \in \{0,1\}} \sum_{\mathbf{x} \in C_i} (\mathbf{x} - \boldsymbol{\mu}_i)(\mathbf{x} - \boldsymbol{\mu}_i)^\top$, the Fisher Discriminant Ratio is
\begin{equation}
    \text{FDR} = (\boldsymbol{\mu}_0 - \boldsymbol{\mu}_1)^\top \mathbf{S}_W^{-1} (\boldsymbol{\mu}_0 - \boldsymbol{\mu}_1).
    \label{eq:fdr}
\end{equation}
This quadratic form summarizes separation relative to within-class scatter \cite{duda_pattern_2001}. When $\mathbf{S}_W$ is rank deficient, the value depends on the regularization convention. Chapter~4 uses $(\mathbf{S}_W+\epsilon\mathbf{I})^{-1}$ with $\epsilon=10^{-4}\operatorname{tr}(\mathbf{S}_W)/d$. Its large full-sample value is therefore a descriptive, scale- and regularization-dependent statistic rather than an out-of-sample effect size.

\subsection{The Information Bottleneck Principle}
\label{sec:info_bottleneck}

The representation pipeline developed in Chapters~\ref{chap:lsm}--\ref{chap:embeddings} can be understood through the lens of the \emph{information bottleneck} (IB) principle \cite{tishby_deep_2015}. Given an input variable $X$ (the EEG signal), a target variable $Y$ (the affective condition or diagnostic label), and a compressed representation $T$ (the ARSPI-Net embedding), the IB framework seeks a representation that minimizes the mutual information $I(X; T)$ (compression) while maximizing the mutual information $I(T; Y)$ (relevance):
\begin{equation}
    \min_{p(t|x)} \; I(X; T) - \beta \, I(T; Y)
    \label{eq:ib}
\end{equation}
where $\beta$ controls the compression--relevance trade-off. Dimensional expansion or compression does not by itself determine mutual information. This dissertation does not estimate $I(X;T)$ or optimize the IB objective, so the framework is used only as background vocabulary rather than as an empirical explanation of the pipeline.

The better archived accuracy of BSC$_6$ than mean firing rate shows that the chosen temporal code is more useful to the tested decoder. The panel-centering result shows that subtracting the within-subject mean changes the task geometry. Neither result quantifies a movement along an information-bottleneck curve or proves selective removal of subject information.

\subsection{Elements of Graph Theory and Graph Signal Processing}
\label{sec:graph_theory}

The graph-based component of ARSPI-Net (Chapter~5) relies on several constructs from algebraic graph theory and graph signal processing.

\paragraph{Graph Definitions.}
A graph $G = (V, E)$ consists of a set of $N = |V|$ nodes and a set of edges $E \subseteq V \times V$. For weighted graphs, an edge $(i,j)$ carries a weight $w_{ij} \geq 0$. The graph structure is encoded in the adjacency matrix $\mathbf{A} \in \mathbb{R}^{N \times N}$, where $A_{ij} = w_{ij}$ if $(i,j) \in E$ and $A_{ij} = 0$ otherwise. The degree of node $i$ is $d_i = \sum_j A_{ij}$, and the degree matrix $\mathbf{D} = \text{diag}(d_1, \ldots, d_N)$ collects these on the diagonal.

\paragraph{The Graph Laplacian.}
The combinatorial graph Laplacian is defined as $\mathbf{L} = \mathbf{D} - \mathbf{A}$. Its normalized form, used in spectral graph convolutions, is
\begin{equation}
    \mathbf{L}_{\text{norm}} = \mathbf{I}_N - \mathbf{D}^{-1/2} \mathbf{A}\, \mathbf{D}^{-1/2}.
    \label{eq:normalized_laplacian}
\end{equation}
Both forms are symmetric positive semidefinite. Their eigendecomposition can be written
\begin{equation}
\mathbf{L}=\mathbf{U}\boldsymbol{\Lambda}\mathbf{U}^{\top}, \qquad
\mathbf{U}=[\mathbf{u}_1,\ldots,\mathbf{u}_N], \qquad
\boldsymbol{\Lambda}=\operatorname{diag}(\lambda_1,\ldots,\lambda_N).
\end{equation}
The columns of $\mathbf{U}$ are orthonormal, and $0=\lambda_1\leq\cdots\leq\lambda_N$. The eigenvectors form the graph Fourier basis. Small eigenvalues correspond to signals that vary slowly across strong graph edges, while large eigenvalues correspond to faster graph variation \cite{shuman_emerging_2013}.

\paragraph{Spectral Graph Convolution.}
For a signal $\mathbf{x} \in \mathbb{R}^N$ defined on the nodes of $G$, a spectral graph convolution with a filter $g_\theta$ is defined as
\begin{equation}
    g_\theta \star \mathbf{x} = \mathbf{U}\, g_\theta(\boldsymbol{\Lambda})\, \mathbf{U}^\top \mathbf{x},
    \label{eq:spectral_conv}
\end{equation}
where $g_\theta(\boldsymbol{\Lambda}) = \text{diag}(g_\theta(\lambda_1), \ldots, g_\theta(\lambda_N))$ is a learned spectral filter. Direct computation of Equation~\eqref{eq:spectral_conv} requires eigendecomposition of $\mathbf{L}$, which costs $O(N^3)$. Practical GNN architectures approximate $g_\theta$ with low-order polynomials of $\mathbf{L}$ to avoid this cost: ChebyNet uses Chebyshev polynomials \cite{Defferrard2016}, and the Graph Convolutional Network (GCN) uses a first-order approximation that yields the propagation rule $\mathbf{H}^{(k+1)} = \sigma(\tilde{\mathbf{D}}^{-1/2} \tilde{\mathbf{A}}\, \tilde{\mathbf{D}}^{-1/2} \mathbf{H}^{(k)} \mathbf{W}^{(k)})$, where $\tilde{\mathbf{A}} = \mathbf{A} + \mathbf{I}_N$ includes self-loops \cite{Kipf2017}.

\paragraph{Message Passing.}
Most GNN architectures used in this dissertation can be described within the message-passing neural network (MPNN) framework \cite{Gilmer2017}. At each layer $k$, a node $v$ updates its representation by aggregating messages from its neighbors:
\begin{equation}
    \mathbf{h}_v^{(k+1)} = \psi^{(k)}\!\left(\mathbf{h}_v^{(k)},\; \bigoplus_{u \in \mathcal{N}(v)} \phi^{(k)}(\mathbf{h}_u^{(k)}, \mathbf{h}_v^{(k)}, \mathbf{e}_{uv})\right),
    \label{eq:mpnn}
\end{equation}
where $\phi^{(k)}$ is the message function, $\psi^{(k)}$ is the update function, $\mathbf{e}_{uv}$ denotes edge features, and $\bigoplus$ is a permutation-invariant aggregation operator (e.g., sum, mean, or max). Graph Attention Networks (GATs) instantiate this framework by learning attention coefficients $\alpha_{uv}$ that weight the messages from different neighbors \cite{Velickovic2018}. Chapter~5 develops the specific instantiations used in ARSPI-Net.

\paragraph{Graph-Level Readout.}
For graph classification, predicting a label for an entire graph rather than individual nodes, a readout function aggregates node representations into a single graph-level embedding:
\begin{equation}
    \mathbf{h}_G = \text{READOUT}\!\left(\left\{\mathbf{h}_v^{(K)} \mid v \in V\right\}\right),
    \label{eq:readout}
\end{equation}
where common choices include global mean pooling, sum pooling, or attention-based pooling. In the SHAPE application, each graph represents one trial-averaged participant--condition matrix, nodes are unlabeled channel indices, and the readout produces an observation-level label.

\subsection{Dynamical Systems Concepts}
\label{sec:dynamical_systems}

The dynamical summaries in Chapter~6 require a distinction between autonomous and input-driven dynamics. This subsection introduces the relevant notation.

\paragraph{State Space and Trajectories.}
A discrete-time dynamical system is specified by a state vector $\mathbf{x}[t] \in \mathbb{R}^n$ and an evolution rule $\mathbf{x}[t+1] = \mathbf{F}(\mathbf{x}[t])$. The sequence of states $\{\mathbf{x}[0], \mathbf{x}[1], \ldots\}$ is a trajectory in the $n$-dimensional state space. For the LSM reservoir in this dissertation, the state is the membrane potential vector $\mathbf{M}[t] \in \mathbb{R}^{N_{\text{res}}}$, and the evolution rule includes the input-dependent current and the threshold-and-reset nonlinearity defined in Equation~\eqref{eq:lif_discrete}.

\paragraph{Stability and the Lyapunov Exponent.}
The sensitivity of a dynamical system to small perturbations is quantified by Lyapunov exponents. For a trajectory starting at $\mathbf{x}[0]$, the largest Lyapunov exponent measures the average exponential rate of divergence or convergence of nearby trajectories:
\begin{equation}
    \lambda_{\max} = \lim_{T \to \infty} \frac{1}{T} \sum_{t=0}^{T-1} \ln \frac{\|\delta \mathbf{x}[t+1]\|}{\|\delta \mathbf{x}[t]\|}
    \label{eq:lyapunov}
\end{equation}
where $\delta \mathbf{x}[t]$ is an infinitesimal perturbation evolved through the linearized dynamics \cite{strogatz_nonlinear_2015}. If $\lambda_{\max} < 0$, perturbations shrink exponentially and the system is \emph{contracting}: nearby trajectories converge. If $\lambda_{\max} > 0$, perturbations grow exponentially and the system exhibits \emph{sensitive dependence on initial conditions}, the hallmark of chaos. The boundary $\lambda_{\max} \approx 0$ is the \emph{edge of chaos}, where the system is marginally stable.

\paragraph{Driven Versus Autonomous Dynamics.}
A critical distinction for reservoir computing is between the \emph{autonomous} Lyapunov exponent (computed for the free-running system $\mathbf{x}[t+1] = \mathbf{F}(\mathbf{x}[t])$ with no external input) and the \emph{driven} or \emph{conditional} Lyapunov exponent (computed under a specific input drive $\mathbf{u}[t]$). For the LSM, the relevant quantity is the driven exponent:
\begin{equation}
    \lambda_1^{(\text{driven})} = \lim_{T \to \infty} \frac{1}{T} \sum_{t=0}^{T-1} \ln \frac{\|\mathbf{M}'[t+1] - \mathbf{M}[t+1]\|}{\|\mathbf{M}'[t] - \mathbf{M}[t]\|}
    \label{eq:driven_lyapunov}
\end{equation}
where $\mathbf{M}[t]$ and $\mathbf{M}'[t]$ are trajectories driven by the \emph{same} input but initialized differently. A negative finite-time estimate is evidence of contraction for the tested input and perturbation protocol. It is not, by itself, a proof of the Echo State Property over an input class. Simple spectral-radius heuristics from smooth reservoirs do not supply such a proof for this thresholded implementation \cite{yildiz_re-visiting_2012}.

\paragraph{Practical Estimation.}
For finite time series, $\lambda_{\max}$ is estimated using the Benettin renormalization algorithm \cite{benettin_lyapunov_1980}: a perturbed trajectory is evolved alongside the reference trajectory, and at regular intervals the perturbation vector is rescaled to unit norm while accumulating the logarithmic stretching factors. For spiking networks, the threshold-and-reset discontinuity requires care: when one trajectory crosses threshold while the other does not, a finite state-space jump occurs that is not captured by tangent-space linearization. Chapter~6 addresses this implementation detail in the context of the specific LIF reservoir used.

\paragraph{Complexity Measures.}
Two additional complexity measures are used in Chapter~6. \emph{Lempel--Ziv complexity} counts the growth of distinct parsed substrings under a specified parsing convention and can be normalized for sequence length \cite{lempel_complexity_1976}. \emph{Permutation entropy} summarizes the distribution of ordinal patterns in delay-embedded segments and is invariant to strictly monotone amplitude transformations \cite{bandt_permutation_2002}. Both are algorithm-dependent summaries, not direct measures of biological complexity.

\paragraph{Koopman Operator Theory and Linearization of Nonlinear Dynamics.}
For a discrete-time system $\mathbf{x}[t+1]=\mathbf{F}(\mathbf{x}[t])$, the Koopman operator acts linearly on scalar observables by composition, $(\mathcal{K}g)(\mathbf{x})=g(\mathbf{F}(\mathbf{x}))$, even when $\mathbf{F}$ is nonlinear. Finite-dimensional linear representations require an invariant observable subspace, which need not exist for a chosen dictionary \cite{gulina_koopman_2021}.

The reservoir can be described informally as a nonlinear lift into a higher-dimensional state. However, no Koopman eigenfunctions, invariant subspace, or approximation error are estimated in this dissertation. Linear-readout performance therefore cannot be presented as a Koopman prediction or proof; Koopman theory is only an analogy that motivates future analysis.


%% ===================================================================
%% SECTION 4: RESERVOIR COMPUTING
%% ===================================================================
\section{Reservoir Computing: A Framework for Harnessing Complex Dynamics}
\label{sec:reservoir_computing}

\subsection{Theoretical Underpinnings}
\label{sec:rc_theory}
Reservoir computing fixes the recurrent transformation and trains a readout \cite{jaeger2001echo}. The fixed recurrent weights avoid backpropagating through those dynamics, although readout selection, hyperparameters, stability, and computational cost still require evaluation.

The Echo State Property requires asymptotic loss of initial-condition dependence for inputs in a stated class \cite{jaeger2001echo,yildiz_re-visiting_2012}. Fading memory is related but additionally concerns continuity with respect to recent versus remote input history. Spectral radius below one is a common design heuristic; general sufficient conditions depend on stronger contraction bounds, and neither that heuristic nor a single finite-time Lyapunov estimate proves the property for this LIF reservoir.

Approximation results show that reservoir families can represent broad classes of causal fading-memory operators under stated assumptions \cite{boyd_chua_1985,Maass2002RealTime,gonon_ortega_2020}. These are existence results for suitable systems and readouts. They do not establish separation, approximation accuracy, or universality for the finite random reservoir used in this dissertation.

Dambre et al.\ formalize information-processing capacity and show constraints on how capacity is distributed across tasks \cite{dambre_2012}. Other work studies ordered and critical regimes in recurrent systems \cite{busing_2010}. The dissertation does not locate a phase transition or optimize a theoretical capacity measure; its operating point is selected from a finite empirical grid.

\subsection{Echo State Networks vs.\ Liquid State Machines}
Echo State Networks commonly use continuous-valued units \cite{jaeger2001echo}, whereas Liquid State Machines use spiking units \cite{Maass2002RealTime}. Either can be implemented on conventional hardware, and either may be advanced on a simulation clock. This dissertation uses a discrete-time spiking reservoir because it exposes spike events and membrane trajectories for the planned analyses, not because it proves greater biological realism or hardware efficiency.


%% ===================================================================
%% SECTION 5: AFFECTIVE NEUROSCIENCE
%% ===================================================================
\section{Affective Neuroscience: Task Formulation and Ontological Considerations}
\label{sec:affect_ontology}

The dissertation uses experimenter-assigned affective stimulus conditions as labels. Those categories are discussed with reference to valence and arousal models, but participant- and stimulus-level valence/arousal scores are not modeled.

The valence--arousal circumplex model represents affective states in a continuous two-dimensional space \cite{russell_1980}. It provides context for affective tasks, but the present analyses do not use valence or arousal scores. Other accounts emphasize construction from domain-general processes \cite{barrett_2017} or appraisal of a stimulus's significance \cite{scherer_2009}. These theories make clear that an experimenter-assigned image category is not equivalent to a participant's subjective emotional state.

The SHAPE labels identify stimulus conditions rather than each participant's experienced emotion. A correct condition classifier would therefore predict the experimental category, not prove recovery of subjective affect, valence, or arousal.

This dissertation does not resolve these ontological debates. It evaluates a signal-processing pipeline against the available condition labels and screens associations with overlapping diagnostic metadata. It does not establish subjective emotion decoding, diagnostic biomarkers, or a unitary neurobiological mechanism.


%% ===================================================================
%% SECTION 6: EEG PREPROCESSING
%% ===================================================================
\section{Standard EEG Preprocessing}
\label{sec:preprocessing}

Before any computational model is applied to scalp EEG, the raw signal must be processed through a sequence of standard steps designed to attenuate the artifact contamination and measurement confounds described in Chapter~1. This section briefly reviews the standard pipeline so that the preprocessing choices made in Chapter~5 can be evaluated in context.

The typical pipeline proceeds as follows. First, the continuous EEG is \emph{bandpass filtered} to retain frequencies of interest (commonly 0.1--30\,Hz or 0.1--40\,Hz for ERP studies, wider for oscillatory analyses) and to suppress low-frequency drift and high-frequency muscular contamination. Second, the data is \emph{re-referenced} to a common reference scheme (typically the average of all electrodes, linked mastoids, or the Reference Electrode Standardization Technique (REST) proposed by Yao \cite{yao_2001}) because the choice of reference affects all channel-to-channel correlation estimates and therefore directly impacts graph construction in GNN-based analyses. Third, the continuous data is segmented into \emph{epochs} time-locked to stimulus events, with each epoch spanning a defined pre-stimulus baseline and post-stimulus analysis window.

Fourth, \emph{artifact rejection} is applied to remove or correct epochs contaminated by ocular, muscular, or other non-neural activity. Simple threshold-based rejection discards epochs in which any channel exceeds a voltage criterion (e.g., $\pm 100\;\mu$V). More sophisticated approaches use Independent Component Analysis (ICA) \cite{delorme_eeglab_2004} to decompose the multichannel signal into statistically independent components, identify components corresponding to blinks, saccades, or EMG on the basis of their topography and time course, and subtract them before reconstructing the cleaned signal. ICA-based correction preserves more trials than outright rejection but introduces its own assumptions (statistical independence of sources, linearity of mixing) that may not hold exactly.

Finally, epochs may be \emph{baseline-corrected} by subtracting the mean voltage during the pre-stimulus interval from each time point, removing tonic voltage offsets that are unrelated to stimulus-evoked activity. The specific preprocessing choices adopted in this dissertation are detailed in Chapter~5.


%% ===================================================================
%% SECTION 6: INTERPRETABILITY IN CLINICAL MACHINE LEARNING
%% ===================================================================
\section{Interpretability in Clinical Machine Learning}
\label{sec:interpretability_background}

Any clinical use of an EEG model would require evidence for predictive validity, explanation fidelity, usability, and safety. No clinician study is conducted here. This section therefore uses \emph{interpretability} to mean specific forms of model inspection and traceability, not clinician trust or actionability.

\subsection{The Clinical Requirement for Interpretable AI}

If an EEG model were proposed for clinical use, prediction alone would not establish safety, efficacy, or a physiological explanation. Its intended users would also need evidence about data provenance, uncertainty, failure modes, and the fidelity and usability of any explanation. This dissertation does not evaluate those clinical requirements; it uses traceability and model inspection as narrower engineering goals.

\subsection{Post-Hoc Versus Intrinsic Interpretability}

Interpretability is not a single binary property. It depends on the question, the user, and the evidence that an explanation is faithful and useful. Two broad families of methods are relevant here.

\textbf{Post-hoc methods} attribute a fitted model's output to inputs or intermediate variables. Examples include SHAP \cite{lundberg_shap_2017}, Integrated Gradients \cite{sundararajan_ig_2017}, and gradient saliency. These methods explain aspects of the fitted function under their assumptions; they do not reveal the data-generating process. Attention weights are likewise model variables and are not automatically faithful feature attributions \cite{jain_attention_2019}.

\textbf{Interpretable-by-design models} constrain the computation or expose understandable intermediate quantities. Even then, naming a variable does not establish that it is biologically meaningful, causally faithful, or useful to clinicians. Those properties require separate validation \cite{rudin_stop_2019}.

\subsection{Inspectable Quantities in ARSPI-Net}

ARSPI-Net is designed to expose intermediate quantities. Their supported interpretations and limits are:
\begin{itemize}
    \item \textbf{Reservoir stage}: membrane trajectories and spike rasters are repeatable for a saved input, parameter set, and random initialization. They describe the implemented model, not biological neurons.
    \item \textbf{Temporal coding stage}: BSC$_6$ assigns spike counts to six known early windows. A bin can be traced to a time interval, but its value does not by itself identify an ERP component or explain a prediction.
    \item \textbf{Compression stage}: PCA coordinates are data-dependent linear combinations. Loadings can be inspected, but the archived per-channel bases are not mutually commensurate and were fitted globally in the SHAPE analysis.
    \item \textbf{Dynamical-summary stage}: spike count, firing rate, entropy, variance, sparsity, and autocorrelation summaries have defined computational scales. Only some have units; none is a direct clinical measurement.
    \item \textbf{Graph-summary stage}: graph metrics describe a constructed feature-similarity graph. With unlabeled channels and the archived PCA limitations, they have no anatomical-connectivity interpretation.
    \item \textbf{Cross-summary stage}: correlations between dynamical and graph summaries measure association between two model-derived quantities. They are not validated brain structure--function coupling.
\end{itemize}

The staged signal chain improves auditability because each transformation can be inspected. It does not make every output an explanation, and Chapters~\ref{chap:lsm}--\ref{chap:synthesis} evaluate only selected forms of traceability and association.

The model comparisons in this dissertation differ in architecture, training, and explanation method. Their accuracy gaps cannot be interpreted as the measured cost of interpretability, and no clinician-utility study was conducted for ARSPI-Net.

\subsection{Related Work on Interpretable EEG Classification}
\label{sec:interp_related_work}

The development of interpretable methods for EEG analysis spans three distinct research traditions: post-hoc attribution applied to conventional deep learning, brain-structured spiking neural networks, and prototype-based classification.

\paragraph{Post-hoc attribution for EEG.} Gradient saliency, SHAP, and related methods can be applied to EEG models to identify model-sensitive inputs \cite{shrikumar_deeplift_2017,lundberg_shap_2017}. Bitar et al.\ \cite{bitar_snn_grad_2023} proposed gradient-based attribution for spike-form data, Kim and Panda \cite{kim_sam_2021} proposed Spike Activation Maps, and Yao et al.\ \cite{yao_attention_snn_2023} studied attention in spiking networks. Each method requires fidelity evaluation in the intended task; this review does not claim exhaustive priority or absence of other clinical applications.

\paragraph{Brain-structured spiking networks.} NeuCube \cite{kasabov_neucube_2014} uses supplied brain coordinates and STDP to organize a three-dimensional reservoir. Doborjeh et al.\ \cite{doborjeh_explainable_2021} related learned spatiotemporal patterns to ERP stages, and a later study compared NeuCube with BiLSTM on an epilepsy--migraine task \cite{doborjeh_epilepsy_2024}. ARSPI-Net instead uses fixed random reservoir weights and unlabeled SHAPE channel indices. Fixed weights improve computational repeatability but do not provide anatomical grounding.

\paragraph{Prototype-based classification.} ProtoPMed-EEG \cite{barnett_protopmed_2024} evaluated prototype-based case reasoning for ICU EEG and included a clinician study. It motivates the use of inspectable exemplars. ARSPI-Net's prototypes are learned in its own feature space, but this dissertation does not evaluate whether clinicians understand or benefit from them.

\paragraph{Clinical and regulatory scope.} Transparency can be relevant to medical-device development, but this dissertation is not a medical-device submission and does not test safety, efficacy, human factors, or regulatory conformity. No regulatory or clinician-trust claim is used to validate the architecture.

\subsection{The Accuracy--Interpretability Frontier}
\label{sec:accuracy_interpretability}

Accuracy and interpretability need not form a universal trade-off. Rudin argues for interpretable models in high-stakes settings when they can achieve the required performance \cite{rudin_stop_2019}, and Barnett et al.\ evaluate a prototype approach in a particular ICU EEG task \cite{barnett_protopmed_2024}. These results do not define a task-independent frontier.

The ARSPI-Net and EEGNet experiments use different architectures and attribution mechanisms. Their displayed accuracy difference is not the cost of interpretability. ARSPI-Net's attention profiles do not differ by condition under the reported permutation test ($p=0.634$), and the selected early bins do not identify P300, LPP, or N2 components. No screening or clinical-characterization use is supported by this study.

\subsection{A Four-Level Interpretability Taxonomy for Neuromorphic EEG}
\label{sec:interp_taxonomy}

For this staged pipeline, four questions organize the interpretability analysis. They are an audit framework, not cumulative achievement levels or established clinical standards.

\textbf{Question~1: Temporal traceability.} Can a feature be mapped to a predefined input interval? BSC$_6$ supplies such indexing, but the reported correlations do not validate ERP-component recovery because the early bins do not cover the stated LPP window and only overlap part of the P300 window.

\textbf{Question~2: Representation geometry.} Can sample variation and model sensitivity be summarized? Chapter~\ref{chap:arspi_net} reports decompositions and attention weights, but global PCA, transductive centering, and a nonsignificant attention-profile test limit their interpretation.

\textbf{Question~3: Dynamical characterization.} Can reservoir trajectories be summarized by named computational descriptors? Chapter~\ref{chap:dynamics} defines such descriptors. Their physiological and patient-level meanings are not validated, and several are dimensionless rather than physical units.

\textbf{Question~4: Cross-summary correspondence.} Are dynamical and graph summaries associated? Chapter~\ref{chap:synthesis} reports correlations, but the graph construction and permutation null limit them to descriptive model-to-model correspondence. No added information or clinician-facing consistency measure is established.

The questions are not cumulative: success on one does not validate another. Table~\ref{tab:taxonomy_metrics} summarizes the measurements and their evidentiary status after the post-defense audit.

\begin{table}[t]
\centering
\caption{Audit status of four interpretability questions. The displayed numerical thresholds were not established as preregistered clinical validation criteria.}
\label{tab:taxonomy_metrics}
\small
\begin{tabular}{p{2.5cm}p{3.5cm}p{3.5cm}p{3.5cm}}
\toprule
\textbf{Question} & \textbf{Measurement} & \textbf{Observed result} & \textbf{Audit status} \\
\midrule
Q1: Temporal traceability & Bin-to-time mapping; exploratory ERP-scalar correlations & Early bins map to $39$--$273$~ms; spike-domain $|r|\leq0.33$ & Time indexing supported; ERP recovery not established \\
\midrule
Q2: Representation geometry & Variance summaries; panel centering; attention & Archived $\rho=7.2$; attention permutation $p=0.634$ & Descriptive; global PCA and transductive estimand \\
\midrule
Q3: Dynamical summaries & Seven predefined reservoir summaries & Defined computational scales; small exploratory condition effects & Inspectable model summaries; physiology not validated \\
\midrule
Q4: Cross-summary correspondence & Within-observation cross-channel correlation & Archived median $\kappa=0.273$ & Descriptive; graph basis and null-model limitations \\
\bottomrule
\end{tabular}
\end{table}

The framework helps organize what this implementation exposes and what remains unvalidated. It is not used to rank ARSPI-Net, ProtoPMed-EEG, NeuCube, or EEGNet on a universal interpretability scale, and no priority claim is made.


%% ===================================================================
%% SECTION 7: STATE OF THE ART
%% ===================================================================
\section{Related Work and Positioning}
\label{sec:state_of_art}

\subsection{LSMs and Reservoir Computing in Physiological Signal Processing}
\label{sec:lsm_literature}
Reservoir computing has been applied to several temporal classification problems. Verstraeten et al.\ evaluated reservoir methods on isolated spoken digits \cite{verstraeten_experimental_2007}, and Bozhkov et al.\ used Echo State Networks with intrinsic plasticity for cross-subject EEG valence discrimination \cite{bozhkov_reservoir_2017}. Separately, Montoya-Mart\'{i}nez et al.\ studied the number and placement of electrodes for reconstructing attended-speech envelopes from EEG \cite{montoya-martinez_effect_2021}; that work motivates attention to channel selection but is not a reservoir-size study.

Yang et al.\ report a deep spiking approach to seizure detection \cite{yang_neuromorphic_2023}, and Tanaka et al.\ review physical reservoir computing \cite{tanaka_recent_2019}. Cai et al.\ combine spiking and graph convolution for auditory-attention detection \cite{cai_eeg-based_2024}. These studies use different tasks, sensors, preprocessing, and validation schemes and are therefore context rather than direct benchmarks for SHAPE.

Prior work uses a range of inputs, from hand-crafted spectral summaries to more direct time-series representations. The present dissertation examines one particular design: a fixed spiking reservoir driven by minimally summarized waveforms, followed by explicit temporal coding and channel-level analysis. It does not train the temporal and spatial stages jointly and should not be described as an end-to-end learned system.

\subsection{The Emergence of Hybrid Neuromorphic-GNN Architectures}
Spiking and graph methods can be combined when the modeling problem supplies meaningful nodes, edges, and event-derived features. Dy-SIGN is one example for dynamic link prediction \cite{Yoon2024DySIGN}. The present author's earlier paper introduced an initial ARSPI-Net design on a reduced feature set \cite{lane_arspi_2024}; the dissertation extends and critically evaluates that implementation.

\subsection{GNNs in Affective EEG and the Volume Conduction Constraint}
\label{sec:volume_conduction}
GCN and GAT variants have been applied to electrode-level EEG graphs for emotion-recognition tasks \cite{Jia2024GNNEEGSentiment}. Such graphs are modeling constructs, not direct cortical structural maps. Volume conduction, field spread, reference choice, and electrode geometry can all alter channel correlations \cite{nunez_srinivasan_2006,yao_2001}. In this dissertation the limitation is stronger: verified electrode labels and coordinates were unavailable, and some archived PCA bases were noncommensurate. The graph results are therefore restricted to feature relations among channel indices, not functional or anatomical connectivity.

\subsection{Architectural Alternatives}
\label{sec:architectural_comparison}
A natural question is why this dissertation adopts a fixed spiking reservoir coupled with a graph-based readout rather than the most prominent alternative architectures. Three comparisons clarify the design rationale.

End-to-end SNNs often use surrogate gradients for nondifferentiable spike events \cite{neftci_surrogate_2019}; later methods have improved their training on benchmark tasks \cite{zheng_tdbn_2021,fang_plif_2021}. The fixed reservoir used here avoids training recurrent weights and keeps one transformation available for repeated inspection. This is a design choice, not evidence that trained SNNs are unstable or less interpretable in general.

GNNs can operate on conventional EEG summaries such as band power, differential entropy, or ERP amplitudes \cite{Jia2024GNNEEGSentiment}. Binned reservoir spikes offer a different temporal representation. The experiments compare selected variants but do not show that conventional features cannot represent transients or that the reservoir retains all relevant temporal information.

Transformer models provide another option for EEG sequences. Standard self-attention has quadratic sequence-length cost, while efficient variants reduce it. Liu et al.\ study nonstationary time-series forecasting rather than EEG \cite{liu_nonstationary_2022}; their results caution that normalization can remove task-relevant shifts in some settings but do not establish a Transformer limitation for this dataset.

The fixed-reservoir design is the experimental compromise evaluated here. The dissertation does not establish its superiority to trained SNNs, conventional-feature GNNs, or Transformers.

\subsection{Examples of Inspectable Outputs}
\label{sec:interp_landscape}

Table~\ref{tab:interp_comparison} lists the kinds of inspectable outputs reported for four example EEG approaches. The categories are not a validated ranking scale, and absence from a cited implementation does not mean a method is incapable of providing that analysis.

\begin{table}[t]
\centering
\caption{Examples of inspectable outputs in the cited implementations. Entries describe reported analyses, not universal capabilities or clinical validation.}
\label{tab:interp_comparison}
\small
\begin{tabularx}{\textwidth}{p{2.6cm}XXXX}
\toprule
\textbf{Output} & \textbf{ARSPI-Net} & \textbf{ProtoPMed} & \textbf{NeuCube} & \textbf{EEGNet + attribution} \\
\midrule
Temporal view & Early-bin index; attention not condition-separable & Not central to cited model & Coordinate/time-based spike patterns & Post-hoc temporal attribution \\
Case examples & Learned prototypes; no clinician study & Prototype cases with clinician study & Learned spatiotemporal patterns & Not in cited EEGNet implementation \\
Dynamical summaries & Fixed-reservoir computational descriptors & Not central to cited model & STDP/reservoir activity analyses & Hidden-state or input attributions possible \\
Cross-summary analysis & Descriptive correlation of model summaries & Not reported & Not reported in cited study & Not reported in cited study \\
\midrule
Fixed reservoir in cited model & Yes & N/A & No (STDP adaptation) & N/A \\
Clinician evaluation cited & No & Yes & No & No \\
\bottomrule
\end{tabularx}
\end{table}

ProtoPMed-EEG evaluates prototype cases in an ICU EEG setting \cite{barnett_protopmed_2024}; NeuCube supplies coordinate-structured spike analyses \cite{kasabov_neucube_2014,doborjeh_explainable_2021}; and EEGNet can be paired with attribution methods \cite{lawhern_eegnet_2018}. These studies address different tasks and users. The present work should be compared with them conceptually, not by treating the dissertation's four questions as a common validated scorecard.

ARSPI-Net combines a fixed-reservoir trace, temporal bins, prototypes, and model-derived graph summaries in one implementation. Its contribution is the construction and audit of that pipeline; it does not inherit the clinical evidence of other systems or establish priority over all prior architectures.

\subsection{Identifying the Research Gap: The Need for a Principled Embedding Pipeline}
\label{sec:research_gap}
The practical question motivating this dissertation is how to summarize high-dimensional reservoir spike trains for channel-wise classification and graph experiments. Chapter~4 compares several summaries in a controlled synthetic setting, and Chapter~5 applies the selected representation to SHAPE. The post-defense audit identifies global PCA and cross-channel basis problems in the archived real-data analysis, so the contribution is a transparent, correctable engineering study rather than a claim of a fully validated general methodology.
